% ---- pg-597 （书页 585） ----

\noindent 12.31\quad 考虑\hyperref[tab:11.7]{表 11.7} 中的原油数据. 仿照例 11.14 的做法变换数据. 忽略已知的组归属. 用特定目标的软件 MCLUST,

\noindent (a)\quad 允许 12.5 节中的不同的协方差结构, 最多允许有 $K=7$ 个类, 用 BIC 准则选择一个混合模型.

\noindent (b)\quad 将最佳模型的结果与例 11.14 中的分类结果作比较. 注意一个原油类(one classification)是如何对应若干个聚类(clusters)的.

\phantomsection\label{bib:10}
\section*{参考文献}

\begin{enumerate}
\item Abramowitz, M., and I. A. Stegun, eds. \textit{Handbook of Mathematical Functions}. U.S. Department of Commerce, National Bureau of Standards Applied Mathematical Series. 55, 1964.
\item Adriaans, P., and D. Zantinge. \textit{Data Mining}. Harlow, England: Addison-Wesley, 1996.
\item Anderberg, M. R. \textit{Cluster Analysis for Applications}. New York: Academic Press, 1973.
\item Berry, M. J. A., and G. Linoff. \textit{Data Mining Techniques: For Marketing, Sales and Customer Relationship Management} (2nd ed.) (paperback). New York: John Wiley, 2004.
\item Berthold, M., and D. J. Hand. \textit{Intelligent Data Analysis} (2nd ed.). Berlin, Germany: Springer-Verlag, 2003.
\item Celeux, G., and G. Govaert. “Gaussian Parsimonious Clustering Models.” \textit{Pattern Recognition}, \textbf{28} (1995), 781--793.
\item Cormack, R. M. “A Review of Classification (with discussion).” \textit{Journal of the Royal Statistical Society (A)}, \textbf{134}, no. 3 (1971), 321--367.
\item Everitt, B. S., S. Landau and M. Leese. \textit{Cluster Analysis} (4th ed.). London: Hodder Arnold, 2001.
\item Fraley, C., and A. E. Raftery. “Model-Based Clustering, Discriminant Analysis and Density Estimation.” \textit{Journal of the American Statistical Association}, \textbf{97} (2002), 611--631.
\item Gower, J. C. “Some Distance Properties of Latent Root and Vector Methods Used in Multivariate Analysis.” \textit{Biometrika}, \textbf{53} (1966), 325--338.
\item Gower, J. C. “Multivariate Analysis and Multidimensional Geometry.” \textit{The Statistician}, \textbf{17} (1967), 13--25.
\item Gower, J. C., and D. J. Hand. \textit{Biplots}. London: Chapman and Hall, 1996.
\item Greenacre, M. J. “Correspondence Analysis of Square Asymmetric Matrices,” \textit{Applied Statistics}, \textbf{49}, (2000) 297--310.
\item Greenacre, M. J. \textit{Theory and Applications of Correspondence Analysis}. London: Academic Press, 1984.
\item Hand, D., H. Mannila, and P. Smyth. \textit{Principles of Data Mining}. Cambridge, MA: MIT Press, 2001.
\item Hartigan, J. A. \textit{Clustering Algorithms}. New York: John Wiley, 1975.
\item Hastie, T. R., R. Tibshirani and J. Friedman. \textit{The Elements of Statistical Learning: Data Mining, Inference and Prediction}. Berlin, Germany: Springer-Verlag, 2001.
\item Kennedy, R. L., L. Lee, B. Van Roy, C. D. Reed, and R. P. Lippmann. \textit{Solving Data Mining Problems Through Pattern Recognition}. Upper Saddle River, NJ: Prentice-Hall, 1997.
\item Kruskal, J. B. “Multidimensional Scaling by Optimizing Goodness of Fit to a Nonmetric Hypothesis.” \textit{Psychometrika}, \textbf{29}, no. 1 (1964), 1--27.
\item Kruskal, J. B. “Non-metric Multidimensional Scaling: A Numerical Method.” \textit{Psychometrika}, \textbf{29}, no. 1 (1964), 115--129.
\end{enumerate}
